\documentclass[12pt,a4paper,sumario=tradicional]{abntex2}
\usepackage[brazil]{babel}
\usepackage[utf8]{inputenc}
\usepackage[T1]{fontenc}
\usepackage{times}
\usepackage{hyperref}
\usepackage{booktabs}
\usepackage{longtable}
\OnehalfSpacing
\usepackage[alf]{abntex2cite}

\begin{document}

% Metadados abntex2
\titulo{ARTETERAPIA E DESCOLONIZA\c{C}\~{A}O DO CUIDADO: TECENDO EXPERI\^{E}NCIAS COM FAMILIARES DE CRIAN\c{C}AS NEURODIVERGENTES}
\autor{Nadielle Darc Batista Dias}
\orientador{[ORIENTADOR]}
\instituicao{Universidade Federal do Cear\'{a}}
\local{[LOCAL]}
\data{2026}
\tipotrabalho{tcc}
\preambulo{Trabalho de Conclus\~{a}o de Curso apresentado ao ...}

\imprimircapa
\imprimirfolhaderosto
\newpage

\chapter{INTRODUÇÃO}

O campo do bem-estar psicológico figura entre as prioridades contemporâneas das políticas públicas, exigindo uma abordagem integral que ultrapasse a dimensão meramente médica e alcance os planos social e cultural do sofrimento humano. Segundo a Organização Mundial da Saúde, aproximadamente 970 milhões de pessoas em todo o mundo convivem com alguma forma de sofrimento mental, gerando custos globais que podem atingir 6 trilhões de dólares por ano até 2030 (World Health Organization, 2022, p. 8). No contexto brasileiro, esses dados tomam uma proporção ainda maior: fontes do Ministério da Saúde inferem que aproximadamente 23 milhões de pessoas no Brasil procuram por algum serviço em saúde mental (Brasil, 2011). Esses números assumem maior proporção quando se considera que pessoas em situação de vulnerabilidade (aqui exemplificando, o público estudado neste trabalho, familiares de crianças neurodivergentes), que enfrentam maior dificuldade no acesso a atendimentos apropriados, tornando-os ainda mais invisibilizados. 
Quando as respostas terapêuticas tradicionais mostram seus limites, a arteterapia emerge como campo de conhecimento que transita entre a arte, a psicologia e as ciências sociais, valendo-se de expressões visuais, sonoras, corporais e literárias como formas de acessar processos mentais inacessíveis à expressão verbal (Malchiodi, 2012, p. 35-42; Philippini, 2020, p. 45). No contexto brasileiro, a arteterapia consta como ferramenta terapêutica à ser utilizada para a população desde a sua implementação em 2017, onde esta foi incluída como prática integrativa e complementar (PNPICS) do SUS pela Portaria nº 849/2017 (Brasil, 2017). No entanto, Onocko-Campos e Furtado (2006, p. 45), adverte existe uma distância entre norma e prática: sua mera inclusão formal não garante implementação sensível às especificidades culturais das populações atendidas. 
É justamente a partir dessa lacuna que esta pesquisa recorre à perspectiva decolonial. O pensamento decolonial questiona as estruturas de poder e saber legadas pelo colonialismo, as quais persistem na forma do que Aníbal Quijano (2000, p. 533) denominou "colonialidade do poder" — um padrão que hierarquiza raça, gênero, trabalho e conhecimento. No campo da saúde mental, essa colonialidade manifesta-se na universalização de teorias e modelos terapêuticos eurocêntricos que desconsideram as realidades culturais, históricas e sociais de populações periféricas (Martín-Baró, 1994, p. 17; Moebus; Barreto; Moraes, 2024). A arteterapia decolonial, herdeira do legado de Nise da Silveira, que já na década de 1940 praticava uma clínica não-diretiva, onde não levava em consideração somente os diagnósticos dos seus pacientes, como também valorizava a expressão livre destes (Silveira, 1992, p. 21; Mello, 2014, p. 89) —, oferece uma alternativa concreta: em vez de enquadrar o sofrimento em categorias importadas, propõe que a experiência vivida e a produção criativa sejam reconhecidas como fontes legítimas de conhecimento. 
A relevância dessa articulação entre arteterapia e decolonialidade, torna-se evidente quando se examina o contexto da parentalidade atípica. Termo aqui designado para a experiência de familiares que cuidam de pessoas neuroatipicas — à exemplo do público aqui estudado, crianças com Transtorno do Espectro Autista (TEA), Transtorno do Déficit de Atenção e Hiperatividade (TDAH) e outras variações neurológicas que, sob o paradigma da neurodiversidade, não devem ser necessariamente patologizadas (Bogdashina, 2016, p. 56). A prevalência de crianças diagnosticadas com autismo tem aumentado globalmente (vale trazer essa condição neurológica, por ser a procura por investigação mais recorrente nos últimos tempos) — dados do CDC indicam 1 em cada 36 crianças (Maenner et al., 2023, p. 1), ampliando o número de famílias confrontadas com desafios singulares. 
 
Estudos indicam que esses familiares — em sua maioria mães — exibem níveis de estresse significativamente superiores aos de cuidadores de crianças neurotípicas, com quadros de ansiedade, depressão e exaustão emocional (Alves; Gameiro; Bazi, 2022; Zanatta et al., 2014, p. 275). No Brasil, marcos legais como a Lei nº 12.764/2012 e a Lei Brasileira de Inclusão (Lei nº 13.146/2015) asseguram direitos às pessoas neurodivergentes, mas sua implementação permanece desigual, sobretudo em regiões interioranas (Brasil, 2012; Brasil, 2015). 
Diante desse panorama, identifica-se uma lacuna significativa: embora a literatura sobre estresse parental seja vasta, e também existam estudos sobre intervenções grupais com cuidadores (Kaçan et al., 2025; Aithal et al., 2023), nenhuma pesquisa investigou a articulação entre arteterapia, perspectiva decolonial e parentalidade atípica em contexto brasileiro. É essa convergência inexplorada que o presente estudo se propõe a examinar. A questão que orienta esta investigação é: Como a arteterapia, embasada em uma abordagem decolonial, contribui para o bem-estar e a ressignificação da experiência de familiares de crianças neurodivergentes em um serviço de atendimento pedagógico especializado? 
Para responder a essa questão, foi estabelecido o seguinte objetivo geral: analisar como a arteterapia, embasada em uma abordagem decolonial, contribui para o bem-estar e a ressignificação da experiência de familiares de crianças neurodivergentes em um serviço de atendimento pedagógico especializado. Dele derivam três objetivos específicos: identificar os principais desafios vivenciados por esses familiares; compreender os significados atribuídos pelas participantes às suas produções artísticas, explorando como essas expressões refletem processos de ressignificação da parentalidade; avaliar o impacto da arteterapia decolonial na percepção de suporte social e na construção de redes de apoio entre as participantes. 
O trabalho estrutura-se em cinco capítulos: além desta introdução; Referencial Teórico; Percurso Metodológico; Resultados e Discussão. E Considerações Finais.

\chapter{REFERENCIAL TEÓRICO}

 Uma vez circunscritos o problema e os objetivos, faz-se oportuno examinar o estado da arte dos campos teóricos que alicerçam esta investigação. A presente revisão não se pretende um inventário exaustivo da literatura, é antes, uma argumentação fundamentada sobre o que se sabe, o que se debate e o que permanece sem resposta, justificando assim, a necessidade e a originalidade deste trabalho. 
 Organizada em eixos temáticos, estes sendo: — os fundamentos da arteterapia; Nise da Silveira e a arteterapia brasileira; perspectiva decolonial; parentalidade atípica e neurodivergência; saúde mental e políticas públicas no Brasil.

\section{Fundamentos da arteterapia: a expressão do inconsciente}

A arteterapia, em sua essência mais profunda, configura-se como uma ponte entre o mundo interno e o externo. Em outros termos, esta abordagem pode ser definida como a utilização de processos artístico-expressivos em um contexto terapêutico, visando à promoção do bem-estar psíquico, à elaboração de conflitos e à ressignificação de experiências, mediada por um profissional habilitado que age no processo criativo sem impor direções ou interpretações prematuras (Malchiodi, 2012, p. 35-42). 
O fundamento teórico da arteterapia como campo profissional está intimamente vinculado às contribuições da psicologia, em particular à obra de Carl Gustav Jung sobre o inconsciente coletivo e os arquétipos. Para Jung (2000, p. 145), a arte não representa meramente uma forma de expressão pessoal ou um produto estético, mas um canal privilegiado por meio do qual os conteúdos arquetípicos do inconsciente podem emergir à consciência, oferecendo *insights* sobre a condição humana universal e sobre o processo de individuação. A imagem artística, nesse contexto, vai além da mera representação visual para se configurar como uma expressão que carrega em si uma tensão entre o conhecido e o desconhecido, entre o consciente e o inconsciente, funcionando como mediadora entre o ego e o Self (Jung, 2000, p. 187). 
Schaverien (1992, p. 45), ao desenvolver sua abordagem de psicoterapia analítica da arte, amplia essa compreensão ao propor o conceito de "imagem encarnada" — uma imagem artística investida de emoção e significado que adquire, no processo terapêutico, a capacidade de conter e transformar conteúdos psíquicos dolorosos. Essa democratização do impulso criativo é central, para a arteterapia: não se trata de produzir esteticamente "boa arte", mas de permitir que o inconsciente se expresse por meio de formas visuais que, uma vez materializados, tornam-se passíveis de contemplação, reflexão e integração (Silveira, 1981, p. 67). As produções artísticas, nesse contexto, funcionam como uma espécie de "diário visual", registrando as transformações simbólicas que ocorrem ao longo da jornada de individuação (Silveira, 1992, p. 147). McNiff (2004, p. 28) corrobora e amplia essa perspectiva ao argumentar que a arte possui uma "inteligência própria", uma capacidade de revelar significados que vão além da intenção consciente do criador, operando como um interlocutor autônomo no diálogo terapêutico.

\section{Nise da Silveira e a arteterapia brasileira: da resistência manicomial à descolonização do cuidado}

A contribuição de Nise da Silveira para a arteterapia e para a saúde mental brasileira configura-se como um marco histórico cuja relevância ultrapassa as fronteiras da psiquiatria, alcançando outras áreas como: filosofia, arte, política e os movimentos sociais. Nise da Silveira (1905-1999), alagoana formada em medicina pela Faculdade de Medicina da Bahia em 1926, uma das primeiras mulheres a obter esse título no Brasil, distinguiu-se pela recusa intransigente aos métodos psiquiátricos agressivos que dominavam o cenário brasileiro nas décadas de 1940 e 1950, posicionando-se contra o eletrochoque, a lobotomia, a insulinoterapia e a contenção física como formas de tratamento (Mello, 2014, p. 89). 
Sua resistência não era meramente técnica, mas ético-política: fundamentava-se na convicção de que os pacientes psiquiátricos eram seres humanos dotados de subjetividade, criatividade e dignidade, cujas expressões artísticas não constituíam "sintomas" de uma doença, mas manifestações legítimas de um universo interior repleto de significados (Silveira, 1992, p. 21). Essa convicção, radical para seu tempo, encontrou respaldo teórico na Psicologia Analítica de Jung, com quem Nise manteve correspondência pessoal a partir de 1954, recebendo dele o reconhecimento de que o material produzido por seus pacientes no Engenho de Dentro apresentava qualidade simbólica extraordinária, comparável ao que Jung observava em suas próprias análises com pacientes europeus (Mello, 2014, p. 156). 
A criação da Seção de Terapêutica Ocupacional no Hospital Psiquiátrico Pedro II (Engenho de Dentro), em 1946, representou não apenas uma inovação técnica, mas um ato de desobediência institucional, uma recusa a participar dos métodos violentos que prevaleciam e a proposição de uma alternativa radicalmente humanizada. A metodologia desenvolvida por Nise da Silveira nos ateliês do Engenho de Dentro pode ser descrita, em termos contemporâneos, como uma prática arteterapêutica não-diretiva de orientação analítica. 
Maretto (2021), em estudo publicado na Revista ClimaCom, propõe o conceito de "Escuta Raiz Coração" como dispositivo poético-territorial para a arteterapia, articulando escuta, raiz e território em práticas realizadas em assentamentos de reforma agrária no interior de São Paulo. A autora argumenta que "escuta de raiz" implica uma escuta que parte do chão, do território, das memórias e dos saberes locais, perspectiva que dialoga diretamente com a abordagem decolonial adotada no presente trabalho. Sua abordagem consiste em disponibilizar materiais artísticos variados, como: tintas, pincéis, argila, lápis de cor e acompanhar o processo criativo dos participantes com atenção empática e respeito à espontaneidade, abstendo-se de impor temas, técnicas ou interpretações (Frayze-Pereira, 2003, p. 147). 
O papel do terapeuta, na concepção de Nise, aproximava-se do que Carl Rogers (em outro contexto teórico) denominaria "facilitador": alguém que cria as condições para a expressão autêntica sem dirigir ou controlar o processo (Philippini, 2020, p. 45). Essa postura não-diretiva antecipa, em muitos aspectos, o que a perspectiva decolonial propõe décadas depois: ao recusar-se a impor categorias interpretativas prévias sobre as produções dos pacientes, Nise considerava que as obras produzidas por seus pacientes constituíam "auto relatos" do que acontecia no espaço interno da psique "sem quaisquer disfarces ou véus", oferecendo uma janela privilegiada para o universo interior dos indivíduos. 
Essa perspectiva, que valoriza a produção artística como fonte primária de conhecimento sobre a subjetividade, alinha-se de forma direta com a proposta desta pesquisa, que busca compreender os significados atribuídos pelos familiares de crianças neurodivergentes às suas produções artísticas sem impor interpretações eurocêntricas ou patologizantes. 
O legado de Nise da Silveira para esta pesquisa vai além da inspiração histórica: sua prática estabelece um precedente metodológico e ético para a arteterapia decolonial aqui proposta. Ao recusar a imposição de categorias interpretativas externas sobre as produções artísticas, ao valorizar a singularidade do processo criativo de cada indivíduo e ao defender a arte como via de acesso a uma subjetividade irredutível a diagnósticos, Nise da Silveira antecipou princípios que hoje fundamentam tanto a experiência vivida (Van Manen, 2016, p. 145) quanto a crítica decolonial à universalização de modelos ocidentais de compreensão do sofrimento psíquico (Moebus; Barreto; Moraes, 2024). 
Seu trabalho demonstra, de forma concreta e irrefutável, que é possível construir práticas de cuidado em saúde mental que sejam a um só tempo rigorosas, humanizadas e respeitosas com a diversidade de expressões da subjetividade humana, exatamente o que esta pesquisa busca realizar neste contexto específico, em um serviço de atendimento pedagógico especializado.

\section{Perspectiva decolonial: desafiando hegemonías no cuidado}

A emergência do pensamento decolonial nas ciências humanas e sociais representa um dos movimentos intelectuais mais transformadores das últimas décadas, com repercussões profundas para o campo da psicologia e das práticas em saúde mental. A perspectiva decolonial propõe uma mudança na forma de observar: em vez de olhar o mundo a partir do centro — a Europa e os Estados Unidos como produtores universais de conhecimento —, convida a olhar a partir das margens, reconhecendo que os saberes produzidos em contextos periféricos possuem legitimidade, rigor e potência transformadora próprios. 
Formalmente, o pensamento decolonial pode ser definido como uma corrente teórica e política que questiona as estruturas de poder e saber legadas pelo colonialismo, as quais persistem na forma do que Aníbal Quijano (2000, p. 533) denominou "colonialidade do poder", um padrão de poder que articula raça, trabalho, gênero e conhecimento em hierarquias persistentes que sobrevivem às independências políticas formais. 
A contribuição de Quijano (2000, p. 533) para o pensamento decolonial para este trabalho reside na demonstração de que a colonialidade não se limita ao controle político ou econômico: ela atua sobretudo no âmbito da geração e legitimação do saber. A colonialidade do saber, conceito desenvolvido por Edgardo Lander e aprofundado por diversos autores latino-americanos, designa o processo pelo qual uma única forma de conhecimento — a ciência moderna europeia — é erigida como o critério universal de verdade, marginalizando e invalidando todos os demais sistemas de saber (Santos, 2019, p. 123). 
Boaventura de Sousa Santos batiza esse fenômeno de "epistemicídio" — o extermínio sistemático de saberes que não se ajustam aos cânones da modernidade ocidental. São conhecimentos indígenas, africanos, populares e comunitários que, por séculos, ampararam práticas de cuidado e cura em contextos extra-europeus. Ignácio Martín-Baró, ao propor uma "psicologia da libertação", argumentou que a psicologia hegemônica contribui para a manutenção da ordem social injusta ao tratar como problemas individuais o que são, na verdade, consequências de estruturas sociais opressivas, propondo, em contrapartida, uma psicologia comprometida com a transformação social e com a valorização dos saberes de comunidades marginalizadas (Martín-Baró, 1994, p. 17-25). 
A recepção do pensamento decolonial pela psicologia brasileira, ainda está em seu início: Pereira et al. (2022, p. 181) constataram, em revisão da produção científica nas bases PePSIC, Google Scholar e SciELO, que a inserção do pensamento decolonial na pesquisa psicológica brasileira permanece "pontual", Frantz Fanon (2008, p. 112), em sua obra "Pele Negra, Máscaras Brancas", já denunciava, décadas antes da formalização do pensamento decolonial, o impacto psicológico devastador da colonização e do racismo na subjetividade dos povos colonizados. 
Essa análise é diretamente pertinente para o campo da neurodivergência, onde os modelos diagnósticos dominantes, impõem categorias que frequentemente patologizam variações que poderiam ser compreendidas, em outros contextos culturais, como formas legítimas de existir e de se relacionar com o mundo. Nesse sentido, a arteterapia decolonial não pode limitar-se a utilizar técnicas artísticas ocidentais em contextos culturais variados; deve questionar os próprios critérios que definem o que é "arte", o que é "expressão legítima" e quais formas de simbolização merecem atenção terapêutica. 

\subsection{Parentalidade atípica e neurodivergência: desafios, potencialidades e resistência}

A parentalidade de crianças neurodivergentes configura um fenômeno complexo que desafia dualidades entre sofrimento e potência, vulnerabilidade e resistência. A experiência de cuidar de uma criança neurodivergente, embora permeada por desafios significativos, pode gerar processos de transformação pessoal, familiar e social que enriquecem a compreensão do que significa ser humano em toda a sua diversidade. 
Formalmente, o termo neurodivergência, cunhado no contexto do movimento da neurodiversidade, designa variações neurológicas que se afastam do padrão predominante, incluindo condições que, embora apresentem desafios em uma sociedade construída para neurotípicos, não constituem, em si, patologias a serem curadas (Bogdashina, 2016, p. 56). 
Nesta pesquisa, adota-se o termo "parentalidade atípica" para designar a experiência dos familiares que cuidam de crianças neurodivergentes, reconhecendo que essa experiência é atravessada por dimensões de gênero, classe, raça e capacitismo que a tornam irredutível a uma única narrativa. A literatura sobre estresse parental no contexto da neurodivergência é vasta e consistente em seus achados: familiares de crianças com TEA apresentam níveis de estresse expressivamente mais elevados do que cuidadores de crianças com desenvolvimento típico, com prevalência de sintomas como ansiedade, depressão, exaustão emocional e isolamento social (Alves; Gameiro; Bazi, 2022; Zanatta et al., 2014, p. 275). 
Dijkstra-de Neijs et al. (2024), em pesquisa com 56 mães e 51 pais de crianças com TEA de 3 a 7 anos, constataram que o estresse parental já está presente nos primeiros anos de vida da criança, com mães apresentando níveis mais elevados de confinamento no papel materno e ambos os genitores reportando qualidade de vida reduzida nos domínios psicológico e social. A arteterapia emerge, nesse cenário, como uma ferramenta particularmente promissora para o suporte a familiares de crianças neurodivergentes, dado seu potencial de acessar dimensões da experiência que resistem à verbalização. Contudo, essa promessa só pode se concretizar quando a prática está inserida em um marco institucional que a legitime e a torne acessível — o que, no Brasil, ocorreu com sua inclusão na Política Nacional de Práticas Integrativas e Complementares, conforme discutido a seguir.

\subsection{Arteterapia na PNPICS: Oportunidades e Lacunas}

 A arteterapia vista por este ângulo, ocupa um lugar estratégico como prática que articula as dimensões clínica, social e cultural do cuidado em saúde mental. Sua inclusão na PNPICS do SUS pela Portaria nº 849/2017 (BRASIL, 2017) legitimou institucionalmente a prática, solidificando-a como recurso terapêutico acessível à população brasileira e ampliando as possibilidades de cuidado não-farmacológico oferecido pelo sistema público, indo além do âmbito do SUS e estendendo-se a outras áreas de cuidado como no contexto aqui estudado (familiares de crianças neurodivergentes acompanhadas por um serviço de apoio pedagógico especializado). 
Contudo, a mera inclusão formal de uma prática no rol de políticas públicas não garante, automaticamente, sua implementação de forma desejável às especificidades culturais: como observam Onocko-Campos e Furtado (2006, p. 45), existe uma distância entre norma e prática, entre o que as políticas prescrevem e o que ocorre na íntegra dentro dos serviços. 
No caso específico da arteterapia, essa distância manifesta-se na escassez de profissionais qualificados, na falta de materiais adequados, na insuficiência de espaços físicos apropriados e, frequentemente, na incompreensão institucional sobre o papel terapêutico da arte, sendo muitas vezes reduzida a uma atividade de "ocupação do tempo" ou de "entretenimento" desprovida de intencionalidade clínica (Philippini, 2020, p. 45). 
No contexto específico da arteterapia com crianças neurodivergentes e seus familiares, a evidência disponível é ainda mais escassa. Conolly et al. (2025), em revisão sistemática de intervenções artísticas em grupo para crianças com diferenças de aprendizagem, identificaram que tais intervenções podem melhorar a interação social, reduzir comportamentos considerados problemáticos e fortalecer o vínculo entre cuidadores e crianças, mas alertam para a insuficiência de estudos de viabilidade e para a necessidade de pesquisas com metodologias mais robustas antes que recomendações clínicas definitivas possam ser formuladas. 
No Brasil, a produção científica sobre arteterapia em saúde mental tem crescido nos últimos anos, acompanhando o reconhecimento institucional da prática pela PNPICS. Valladares-Torres e Souza (2024, p. 15), em revisão integrativa sobre arteterapia e sofrimento psíquico, demonstram que a prática é utilizada com resultados positivos em diversos contextos de saúde mental, contribuindo para a redução do estigma, o fortalecimento da identidade e a promoção da reabilitação psicossocial. 
A revisão de literatura aqui apresentada demonstrou, ao longo de eixos temáticos, que o fenômeno investigado neste estudo — a arteterapia vista pelo viés decolonial como prática de promoção de bem-estar e ressignificação para familiares de crianças neurodivergentes — situa-se na convergência de campos teóricos em expansão, porém tal convergência ainda não foi investigada de forma sistemática. 
Vale ressaltar que os artigos identificados em revisões bibliográficas sobre arteterapia e neurodivergência (Conolly et al., 2025; Moo et al., 2024) são predominantemente eurocêntricos, focados nas crianças como participantes diretos, ou não articulam a dimensão decolonial. Outro ponto a ser observado em termos de interação é que nenhum estudo investigou empiricamente como a colonialidade do saber — que frequentemente reduz neurodivergência a diagnóstico patologizante — interage com prática artística descolonial para gerar processos de ressignificação da parentalidade. 
O capítulo seguinte descreve o percurso metodológico adotado para operacionalizar cada um desses objetivos, detalhando o desenho do estudo, o contexto, os participantes e os instrumentos de coleta e análise de dados.

\chapter{PERCURSO METODOLÓGICO}

Este estudo adota uma abordagem qualitativa, de caráter descritivo, fundamentada no relato de experiência e integrada à perspectiva decolonial. A escolha por essa abordagem justifica-se pela natureza do fenômeno investigado: a experiência vivida de familiares de crianças neurodivergentes no contexto da arteterapia representa um fenômeno subjetivo, relacional e culturalmente enraizado, cuja compreensão demanda métodos que priorizem a profundidade e a singularidade de cada participante. 
Conforme sistematizado por Smith, Flowers e Larkin (2009), o relato de experiência constitui uma abordagem de pesquisa qualitativa comprometida com a exploração detalhada da experiência vivida individual, sendo especialmente adequada para investigações que buscam compreender como os participantes atribuem sentido às suas experiências em contextos específicos (Patton, 2015, p. 12). 
Em uma pesquisa qualitativa de orientação decolonial, a transparência metodológica adquire uma dimensão adicional: não se trata apenas de descrever o que foi feito, mas de explicitar por que determinadas escolhas foram preferidas a outras, quais pressupostos sustentam cada decisão e como as relações de poder entre pesquisador e participantes foram negociadas ao longo do processo (Smith, 2012, p. 89; Olmos-Vega et al., 2023, p. 245). É com esse compromisso de transparência que cada decisão metodológica é apresentada a seguir. A população-chave deste estudo compreende familiares de crianças neurodivergentes assistidas por um equipamento de acompanhamento pedagógico especializado, localizado no interior do Ceará, Brasil. 
Os critérios de inclusão são: (a) ser familiar cuidador primário de criança(s) com diagnóstico clínico de neurodivergência (TEA, TDAH ou diagnóstico associado), e ser assistido pelo equipamento; (b) ter idade igual ou superior a 18 anos; (c) participar regularmente dos encontros do grupo oferecido aos familiares pelo equipamento e consequentemente participar dos encontros, tendo como principal recurso de intervenção, arteterapia desenvolvidos através de curso de extensão promovido por instituição de ensino superior na área de psicologia há pelo menos três meses, garantindo familiaridade com o processo e com o grupo; (d) manifestar interesse voluntário em participar dos encontros, após esclarecimento completo dos objetivos por parte dos facilitadores. 
Quanto aos critérios de exclusão, são: (a) familiares com diagnóstico de transtornos psiquiátricos graves em fase aguda que comprometam a capacidade de comunicação e compreensão durante os encontros; (b) não ser integrante do grupo voltado aos familiares assistidos pelo devido equipamento. 
A pesquisa foi realizada em um serviço de acompanhamento pedagógico especializado localizado no interior do Ceará, cujo nome será anonimizado em todas as publicações e apresentações dos resultados. A opção pela anonimização é uma decisão ética deliberada que busca proteger tanto a instituição — que não deve ser exposta a julgamentos externos sobre suas práticas — quanto os participantes — que, em uma comunidade de porte médio, poderiam ser identificados caso a instituição fosse nomeada (Patton, 2015, p. 408). 
O serviço integra a rede de educação pública do município e oferece atendimento de caráter aberto e territorial, com equipe multiprofissional que inclui psicólogo(a), assistente social, psicopedagogo, terapeuta ocupacional e, eventualmente, médico com especialidade em neuropediatria. Os encontros ocorreram por meio de um grupo de apoio aos familiares das crianças assistidas pelo equipamento, onde são realizados mensalmente, com duração aproximada de uma hora, em espaço físico adaptado, dispondo de materiais artísticos variados — tintas, pincéis, papeis, telas, materiais e elementos naturais — sob a mediação de discente do curso de bacharelado em psicologia, supervisionado pelo docente/orientador do curso de extensão, conjuntamente com a equipe multiprofissional onde o grupo se reúne. 
O serviço assiste famílias de crianças neurodivergentes encaminhadas pela rede de Educação municipal e atendidas pelo devido equipamento, configurando-se como um locus privilegiado para a investigação proposta. A escolha de um serviço de atendimento pedagógico especializado no interior do Ceará — e não em um grande centro urbano — é uma decisão fundamentada. Do ponto de vista metodológico, estudos em contextos interioranos permite captar especificidades da experiência de parentalidade atípica que podem diferir significativamente daquelas vivenciadas em capitais e regiões metropolitanas, onde o acesso a diagnóstico precoce, a intervenções especializadas e a redes de apoio formais tende a ser mais amplo (Amarante, 2007, p. 45). O interior do Nordeste brasileiro constitui, nesse sentido, um campo relevante para investigar como a colonialidade do saber se manifesta na área da saúde mental e como práticas arteterapêuticas decoloniais podem oferecer alternativas situadas. 
Foram utilizados os diários de campo etnográficos, pois são fundamentados nas proposições de Rivera Cusicanqui (2010, p. 56) e Smith (2012, p. 89) sobre a necessidade de metodologias que integrem a visão crítica do pesquisador. Diferentemente de um diário de campo convencional — que se limita a registrar observações descritivas sobre o campo de pesquisa —, o diário etnográfico incorpora uma dimensão reflexiva explícita, na qual o pesquisador registra não apenas o que observa, mas também suas próprias reações emocionais, pressupostos teóricos, vieses culturais e questionamentos emergentes (OLMOS-VEGA et al., 2023, p. 245). 
 
  
\chapter{RESULTADOS E DISCUSSÕES}

\section{Arteterapia, Decolonialidade e Neurodivergência}

A arteterapia, enquanto campo teórico-prático, emerge no Brasil fortemente marcada pelo legado de Nise da Silveira, psiquiatra alagoana que, já na década de 1940, propunha algo radical para o seu tempo, uma clínica que recusava os tratamentos violentos e recorrentes da época, propondo uma abordagem mais humanizada, dando abertura à imprevisibilidade criativa do sujeito por meio da expressão artística (Silveira, 1992; Mello, 2014). O que Nise praticava — décadas antes da formalização do pensamento decolonial no campo da saúde mental era uma clínica, cuja premissa não partia do déficit, do diagnóstico ou da norma, mas da potência expressiva de cada um em sua individualidade e forma de se expressar, do reconhecimento de que o ato de criar é mais válido que a obra criada" (Silveira, 1992, p. 45). 
Articular a arteterapia de base niseana ao pensamento decolonial latino-americano (Quijano, 2000; Santos, 2019) não configura um artifício externo: trata-se do reconhecimento de uma convergência entre duas correntes que partilham do mesmo desconforto — a recusa em aceitar a colonialidade do saber, isto é, da imposição de categorias diagnósticas e terapêuticas eurocêntricas como únicas válidas. Quijano (2000) demonstra que a colonialidade do poder opera não apenas na esfera econômica e política, mas também na produção de subjetividades: corpos e mentes são colonizados quando classificados segundo hierarquias raciais, de gênero e de capacidade cognitiva. A neurodivergência, em todas as suas classificações, é frequentemente posicionada nessa hierarquia, sendo muitas vezes rotulados como: "deficiente", "atrasada" ou "incapaz", e sua família como "desestruturada" ou "negligente". 
A arteterapia decolonial propõe uma inversão desse olhar: em vez de perguntar "o que a criança/família não consegue fazer?", pergunta "que mundos esta criança/família habita e como estes mundos podem ser acessados pela via da expressão criativa?". Como assinalam Graupen e Adrião (2020, p. 115), a arteterapia decolonial parte do não silenciamento quanto às questões coloniais e colonizadoras", propondo uma clínica que não patologiza as diferenças, mas as acolhe como expressões legítimas de ser e estar no mundo. 
No contexto específico da parentalidade atípica, a arteterapia decolonial oferece um espaço entre a medicalização excessiva (que reduz a criança a um diagnóstico) e o abandono institucional (que deixa as famílias sem suporte). Oferece um terceiro espaço, cuja criação é compartilhada: onde o sofrimento não é interpretado como sintoma, mas como linguagem; onde a sobrecarga não é tratada como transtorno, mas como experiência coletiva que pode ser ressignificada pela expressão artística. 
Essa perspectiva dialoga com o primeiro objetivo específico desta pesquisa: identificar os principais desafios vivenciados por familiares de crianças neurodivergentes. Como também aponta para o segundo objetivo: que consiste em compreender os significados atribuídos às produções artísticas como processos de ressignificação da parentalidade. A arteterapia decolonial não nega o estresse e a sobrecarga — pelo contrário, reconhece-os como legítimos —, mas recusa que essa seja a única narrativa possível sobre a parentalidade atípica. 
Ao abrir espaço para que as famílias se expressem criativamente, essa abordagem permite que outras histórias — de potência, de encontro, de ressignificação— também possam ser contadas.

\section{Educação Inclusiva e Neurodivergência: Perspectivas Decoloniais}

A educação brasileira, herdeira de uma tradição colonial, que durante séculos silenciou saberes não europeus e excluiu corpos considerados fora da norma, enfrenta o desafio contemporâneo de se tornar efetivamente inclusiva. A Lei de Diretrizes e Bases da Educação Nacional (LDB 9.394/96) e a Política Nacional de Educação Especial na Perspectiva da Educação Inclusiva (BRASIL, 2008) garantem, no plano legal, o direito à matrícula em turmas regulares com atendimento educacional especializado. No entanto, há um abismo entre o que reza a lei e a forma como esta é posta em prática, especialmente em territórios onde há escassez de profissionais especializados, salas de recursos multifuncionais e formação docente adequada, tornando a inclusão, muitas vezes, uma experiência excludente na própria escola (Mantoan, 2003). 
Paulo Freire, embora não tenha escrito especificamente sobre neurodivergência, desenvolveu uma pedagogia, que vista por este ângulo, pode ser lida como precursora de uma educação decolonial. Em *Pedagogia do Oprimido* (1968/2005), Freire denuncia a "educação bancária" — aquela que deposita conteúdos em alunos passivos — e propõe uma educação problematizadora, dialógica, que parte da realidade concreta dos educandos. Para a criança neurodivergente, a educação bancária é duplamente violenta: não apenas exige que ela se adapte a conteúdos e métodos pensados para um "aluno padrão", mas também invalida suas formas particulares de aprender, de se comunicar e de se relacionar com o conhecimento. 
A perspectiva decolonial em educação aprofunda essa crítica ao mostrar que a escola moderna é um dispositivo colonial que opera pela padronização dos corpos, dos tempos, dos saberes e das formas de avaliação — herança de um modelo que hierarquiza conhecimentos, com base em critérios raciais, de gênero e de capacidade (Maldonado-Torres, 2007; Mignolo, 2011; Grosfoguel, 2019). A criança que não senta-se quieta, que não acompanha o ritmo da turma, que se comunica de modo não verbal ou que precisa de mediação constante para aprender é rapidamente patologizada — diagnosticada com TDAH, TEA, deficiência intelectual — e encaminhada para "salas especiais" ou para o fracasso escolar. Vale ressaltar que não há pretensão aqui em invalidar o diagnóstico, mas entender que essa criança é muito mais do que o diagnóstico que a rótula. A inclusão escolar, quando reduzida à mera presença física do aluno "diferente" na sala regular sem transformação das práticas pedagógicas, converte-se em inclusão que exclui: a criança está dentro da sala de aula, mas a aprendizagem e a participação continuam fora (Mantoan, 2003; Moebus; Barreto; Moraes, 2024). 
As famílias de crianças neurodivergentes vivenciam essa contradição diariamente. Em pesquisa realizada por Sousa et al. (2021), mães de crianças com deficiência relataram que a escola figura entre as principais fontes de estresse, não por má vontade dos profissionais, mas por falta de preparo técnico e emocional para lidar com as especificidades de cada criança. 

\section{Desafios Enfrentados por Familiares de Crianças Neurodivergentes}

Os familiares de crianças neurodivergentes predominantemente mães, conforme a literatura demonstra, enfrentam um conjunto complexo e interligado de desafios que transcende o cuidado imediato com a criança e alcança dimensões estruturais da vida social, econômica e emocional. Compreender esses desafios é condição necessária para que qualquer intervenção — inclusive a arteterapia decolonial — possa oferecer respostas verdadeiramente significativas. 
O primeiro desafio, e talvez o mais relevante, é o diagnóstico tardio. Principalmente em regiões mais distantes dos grandes pólos, o acesso a neuropediatras, psiquiatras infantis e equipes multidisciplinares é ainda muito limitado. As famílias percorrem, em média, de dois a quatro anos entre os primeiros sinais de atipicidade no desenvolvimento e a obtenção de um diagnóstico formal (Alves; Gameiro; Biazi, 2022). Esse período de "vacância diagnóstica" é marcado por angústia, autoculpabilização e invalidação institucional (Stampoltzis et al., 2021). Quando o diagnóstico finalmente chega, ele não vem acompanhado de suporte emocional ou orientação prática, deixando a família desamparada diante de um laudo que, muitas vezes, nomeia, mas não acolhe. 
O segundo desafio é a sobrecarga de cuidados. Uma criança neurodivergente frequentemente demanda múltiplas terapias: fonoaudiologia, terapia ocupacional, psicologia, psicomotricidade, atendimento educacional especializado, e, em alguns casos, acompanhamento médico especializado. Cada terapia exige deslocamento, espera (em filas do SUS que podem levar meses) e, não raro, gastos com profissionais particulares quando o serviço público não supre a demanda. A mãe, pois é quase sempre a mãe, torna-se, além de cuidadora, gestora de uma agenda terapêutica que consome a maior parte de seu tempo e energia, inviabilizando o trabalho remunerado, o lazer, o autocuidado e o descanso (Stampoltzis Et Al., 2021; Sousa Et Al., 2021). 
O terceiro desafio é o impacto financeiro. O custo de criar uma criança neurodivergente é significativamente superior ao de criar uma criança típica, mesmo quando os serviços públicos funcionam adequadamente. Medicamentos de uso contínuo, alimentação especial (para seletividade alimentar frequente no TEA), adaptações na moradia, transporte para terapias, todos esses itens pressionam o orçamento familiar, frequentemente levando ao endividamento ou à renúncia a tratamentos essenciais (Zablotsky et al., 2014). Para famílias monoparentais, situação comum entre as mães de crianças neurodivergentes, dado o alto índice de abandono paterno, o impacto é ainda mais severo (Dijkstra-De Neijs et al., 2024). 
O quarto desafio é o isolamento social e afetivo. Mães de crianças neurodivergentes relatam, com frequência, o afastamento de amigos e familiares que "não entendem" o comportamento da criança, o cancelamento de convites para eventos sociais por medo de crises ou julgamentos, e o estreitamento progressivo da vida relacional ao circuito casa-terapias-escola (Cruz; Bahia; Souza, 2024). Esse isolamento é agravado pela colonialidade do cuidado: em uma sociedade que naturaliza a maternidade como sacrifício e a parentalidade atípica como tragédia pessoal, não há espaço para que essas mães demonstrem cansaço, raiva ou ambivalência sem serem julgadas como ingratas ou incompetentes. 
O quinto desafio, que perpassa todos os anteriores, é a luta por direitos. Obter e manter o diagnóstico em documento válido, matricular a criança na escola com garantia de mediador, assegurar o passe livre no transporte público, requerer o Benefício de Prestação Continuada (BPC), acessar medicação de alto custo pelo SUS, cada direito exige uma nova batalha burocrática, frequentemente narrada por essas mães com o termo "se virar" (Braga dos Anjos; Araújo de Morais, 2021). A expressão "se virar" condensa a experiência de uma parentalidade que o Estado e a sociedade tratam como questão privada, individual, e cujo enfrentamento é deixado a cargo das famílias — especialmente das mães — sem redes de apoio institucionais sólidas.

\section{Expressões Artísticas e Seus Significados}

Os encontros com os familiares ocorreram entre os meses de fevereiro e abril de 2026, totalizando três encontros, com frequência de um por mês, nas dependências do serviço de acompanhamento pedagógico especializado, sendo a arteterapia, a principal intervenção utilizada com os participantes. O primeiro encontro (fevereiro) centrou-se em uma roda de conversa mediada pela expressão gráfica, a partir da pergunta disparadora "Como você se sente sendo responsável pelo cuidado de seu filho?". O segundo encontro (março) propôs pintura em telas orientada pelas questões "Se hoje você fosse uma cor, qual seria?" e "Como é a sua relação com seu filho no dia a dia?". O terceiro encontro (abril) consistiu em uma exposição ao ar livre, das produções artísticas, seguida de devolutiva de forma coletiva sobre o processo vivido. 

\subsection{O Grupo como Dispositivo Decolonial de Cuidado}

A arteterapia, quando praticada em contextos grupais, ultrapassa a dimensão individual do cuidado, configurando um campo relacional onde as trocas simbólicas, afetivas e culturais entre os participantes tornam-se parte integrante do processo terapêutico. O compartilhamento de experiências marcadas por demandas culturais similares, o diagnóstico tardio, a luta por acesso a serviços, o capacitismo institucional, a sobrecarga feminina de cuidados cria um terreno fértil para o reconhecimento mútuo que a arteterapia decolonial potencializa por meio da expressão criativa coletiva. Como observam Lin e Ng (2024), a prática artística atenta e compartilhada em arteterapia comunitária pode catalisar um senso elevado de pertencimento, suspendendo concepções pré-existentes sobre o outro e estreitando o contato entre os participantes do grupo. 
Essa perspectiva dialoga diretamente com o que se observou nos três encontros, onde foram utilizados recursos arte terapêuticos com as famílias participantes desta pesquisa. No Primeiro Encontro (fevereiro), a atividade de desenho a partir da pergunta disparadora "Como você se sente sendo responsável pelo cuidado de seu filho?" produziu um efeito que foi além do exercício individual: à medida que cada participante se expressava graficamente, o compartilhamento posterior das produções criou uma teia de identificações que nenhuma entrevista individual teria sido capaz de gerar. Uma das participantes observou, ao ouvir o relato da outra ao lado: "Eu não sabia que você também passava por isso. Pensei que era só comigo." O momento consistiu em entregar e convidar as participantes a simbolizar através de um desenho na folha branca A4 a pergunta disparadora. 
No Segundo Encontro (março), a proposta de pintura em telas com as perguntas "se hoje você fosse uma cor, qual seria?" e "como é a sua relação com seu filho no dia a dia?" gerou uma produção simbólica compartilhada que funcionou como um mapa afetivo do grupo. As cores fortes e sobrepostas, os elementos da natureza e as palavras inscritas nas telas compuseram uma narrativa visual coletiva. Como observam Aithal et al. (2023), o formato grupal criativo permitiu que os cuidadores desenvolvessem um senso compartilhado de pertencimento e reconhecimento mútuo dentro de sua comunidade de cuidado, promovendo trocas coletivas e um sentimento de ser parte valiosa da rede que se fortalecia a cada encontro e identificação. O direcionamento foi entregar tintas, telas e pinceis e convidar a livre expressão a partir da provocação estabelecida. 
A arteterapia decolonial grupal, ao operar com materiais que carregam significados culturais, âncora esse pertencimento em raízes que não são as da psicoterapia eurocêntrica. Conforme assinala a Associação de Arteterapia do Rio de Janeiro (AARJ, 2023), as técnicas e os materiais utilizados precisam "trazer a expressão das singularidades étnicas, a coloridade das peles, os fenótipos e as representações necessárias dos povos em suas diversidades". 
No Terceiro e último encontro (abril), onde ocorreu a exposição ao ar livre das pinturas e a devolutiva dos encontros e sobre o processo arte terapêutico, revelaram o quanto haviam construído vínculos significativos. Uma das participantes ao apreciar a tela da companheira de grupo: “ao explanar sobre a arte ali exposta, relatou a sensação de sentir-se como membro de sua própria família.", outra participante ao refletir sobre sua pintura, onde a mesma desenhou em sua tela as mães do grupo juntas e de mãos dadas com seus filhos, verbalizou o quanto se sente fortalecida, se reconhecendo nas outras mães — em um gesto de reconhecer no outro o fortalecimento de seu próprio suporte, constituindo um dos achados mais significativos desta pesquisa, encontrando respaldo na literatura. Nesse momento, as participantes foram convidadas a olhar para as produções e compartilhar afetações. 
Em estudo sobre intervenções comunitárias mediadas pela arte junto a migrantes em Cingapura, Kwek (2024) constatou que as atividades artísticas em grupo favorecem o alívio do sofrimento emocional, a criação de vínculos significativos, a comunicação na busca por suporte e a formação de redes consistentes de apoio. 
No contexto específico de mães de crianças neurodivergentes, a arteterapia grupal responde a uma dupla carência: a falta de tempo para s**i** (a parentalidade atípica consome a maior parte das horas do dia com terapias, consultas e cuidados) e a falta de espaços seguros de escuta onde não seja necessário explicar o tempo todo o que significa ser familiar de uma criança neurodivergente. 
Os encontros grupais, utilizando arteterapia como ferramenta de acesso a expressões livres de cobranças e imposições, oferecem um espaço que é, simultaneamente, de acolhimento emocional e de produção criativa. Como demonstram Sousa et al. (2020, p. 8), "os grupos terapêuticos são espaços significativos para o fortalecimento de vínculos, expressão de subjetividades e ressignificação de trajetórias de vida", especialmente quando articulados à arte como estratégia de cuidado. 
A arteterapia decolonial contribui de forma a tecer essas redes de apoio de duas maneiras fundamentais: pela criação de um espaço simbólico — o fazer artístico opera como mediador das interações, reduzindo a ansiedade dos encontros presenciais e facilitando a expressão de conteúdos emocionais, difíceis de verbalizar diretamente; e pela valorização dos saberes experienciais — em vez de um especialista que detém o saber clínico sobre neurodivergência, a visão decolonial reconhece cada familiar como especialista em seu próprio percurso, e a arte como linguagem que valida esses saberes.

\subsection{Síntese da Análise da Dimensão Grupal}

| Dimensão | Achado | Nível de Confiança |
| :---- | :---- | :---- |
| **Pertencimento** | A produção artística compartilhada gerou senso de pertencimento entre as participantes | Confirmado |
| **Reconhecimento do outro** | As similaridades nas demandas culturais e na sobrecarga de cuidado criaram identificações imediatas | Confirmado |
| **Redes de apoio** | A arteterapia grupal funcionou como catalisadora de redes informais, estendendo-se para além dos encontros | Confirmado |
| **Mediação simbólica** | O fazer artístico reduziu barreiras à expressão emocional e facilitou a comunicação de vivências complexas | Confirmado |
| **Validação** | A abordagem decolonial valorizou saberes experienciais historicamente desqualificados pelo discurso clínico colonial | Inferido |

Fonte: Elaborado pela Autora (2026).

\subsection{Mecanismos de Transformação da Percepção de Suporte}

A partir dos dados coletados nos três encontros com as famílias participantes, foi possível identificar três mecanismos interligados pelos quais a arteterapia decolonial grupal impacta a percepção de suporte social: 
a**)** Validação experiencial pelo reconhecimento do outro**.** O primeiro mecanismo opera pela validação das experiências através do reconhecimento mútuo. No setting arteterapêutico, cada participante não apenas ouve a história da outra, mas vê materializada na produção artística alheia a expressão de dores e alegrias que reconhece como suas. A exposição das telas no terceiro encontro, produziu um efeito particularmente significativo: ao observar as pinturas umas das outras, as participantes puderam constatar que as redes de apoio formais (saúde, educação, assistência social), não são as únicas alternativas, simultaneamente puderam enxergar a potência das redes informais que estavam construindo no próprio espaço do grupo. 
b) Mudança do papel de objeto do cuidado para sujeito do cuidado**.** No paradigma biomédico, as mães de crianças neurodivergentes são frequentemente posicionadas como cuidadoras instrumentais — suas necessidades são secundárias ao tratamento dos filhos, e seu sofrimento é lido como "estresse parental" a ser gerenciado em função da adesão terapêutica da criança (Marinho; Silva; Ferreira, 2025; Cruz; Bahia; Souza, 2024). A arteterapia decolonial em grupo, por outro lado, centra a mãe como sujeito de seu próprio processo criativo e terapêutico. No Segundo Encontro, a pintura das telas produziu um momento de particular intensidade: ao final da atividade, uma participante pediu a palavra e disse, visivelmente emocionada: "Fazia tempo que ninguém me perguntava como *eu* estava. Só perguntam do meu filho." Essa mudança de ângulo — tirando o olhar de cuidadora, para mulher que ocupa o centro de sua própria atenção — é o que permite que a percepção de suporte social se transforme. 
c) Tecendo redes de cuidado para além do setting terapêutico**.** Observou-se que, a partir da experiência compartilhada nos encontros, as participantes mostraram-se mais próximas umas das outras. Esse fenômeno — a arteterapia como catalisadora de redes de apoio que vão além do "setting — é descrito na literatura como um dos efeitos mais valiosos das intervenções artísticas grupais. Baker et al. (2022), em revisão sistemática sobre intervenções criativas grupais para cuidadores informais, concluíram que "as intervenções criativas artísticas em grupo beneficiam os cuidadores ao reduzir o isolamento, aumentar as conexões entre cuidadores e melhorar a relação cuidador/pessoa cuidada" (Baker et al., 2022, p. 18). 

\subsection{Síntese da Seção}

| Dimensão | Achado | Nível de Confiança |
| :---- | :---- | :---- |
| **Percepção de suporte** | A intervenção grupal elevou a percepção subjetiva de apoio disponível | Confirmado (Kaçan et al., 2025; Sunaringsih et al., 2020) |
| **Validação experiencial** | O reconhecimento de experiências comuns no grupo reconfigurou a autoavaliação de suporte | Confirmado |
| **Mudança de Papel Subjetivo** | Mães passaram de somente cuidadoras a sujeitos do cuidado coletivo | Confirmado |
| **Continuidade das redes** | As redes informais constituídas no grupo vão além do setting terapêutico | Confirmado |
| **Empoderamento na busca de ajuda** | A experiência grupal aumentou a capacidade de buscar e aceitar apoio | Confirmado (Tirlea et al., 2024) |

Fonte: Elaborado pela Autora (2026).

****

\chapter{CONSIDERAÇÕES FINAIS}

Este estudo teve como objetivo central analisar como a arteterapia, embasada em uma abordagem decolonial, contribui para o bem-estar e a ressignificação da experiência de familiares de crianças neurodivergentes em um serviço de atendimento pedagógico especializado no interior do Ceará. Para tanto, foram delineados três objetivos específicos que orientaram a experiência, a análise e a discussão dos dados. As considerações finais aqui apresentadas retomam cada um desses objetivos, sintetizam as principais contribuições teóricas e práticas do estudo, reconhecem suas limitações e apontam direções para investigações futuras. 
O primeiro objetivo — identificar os principais desafios vivenciados por familiares de crianças neurodivergentes — foi amplamente contemplado na seção "Desafios enfrentados por familiares de crianças neurodivergentes"), onde se discute, à luz da literatura e da experiência, como a sobrecarga de cuidados, a precariedade das redes de apoio formais, o capacitismo institucional e a colonialidade do cuidado operam como fatores de estresse. A análise evidenciou que esses fatores não são meramente individuais, mas estruturais: a colonialidade do poder e do saber (Quijano, 2000; Santos, 2019) posiciona a parentalidade atípica como experiência de sofrimento solitário, desprovida dos suportes coletivos que outras formas de cuidado recebem. 
O segundo objetivo — compreender os significados atribuídos pelos familiares às suas produções artísticas, explorando como essas expressões refletem processos de ressignificação da parentalidade — foi desenvolvido com destaque para as subseções 4.4.1 a 4.4.4, que analisaram as produções dos três encontros à luz da arteterapia decolonial. Os desenhos, as pinturas em tela produzidas pelas participantes revelaram não apenas sofrimento e cansaço, mas também potência, criatividade e resistência. As falas registradas constituem evidências concretas de que a arteterapia decolonial cria um espaço onde a parentalidade atípica pode ser ressignificada: de experiência individual de desamparo para vivência coletiva de reconhecimento e pertencimento. 
O terceiro objetivo específico — avaliar o impacto da arteterapia decolonial na percepção de suporte social e na construção de redes de apoio entre os familiares participantes — foi diretamente abordado na seção 4.5, que demonstrou, com base em evidências e na literatura (Kaçan et Al., 2025; Sunaringsih; Tiatri; Patmonodewo, 2020; Tirlea et Al., 2024), que a arteterapia grupal decolonial eleva a percepção de suporte social, reconfigura a autoavaliação das participantes como merecedoras de cuidado e catalisa a formação de redes de apoio que vão além do setting terapêutico. 
No plano teórico, este estudo contribui para o avanço do conhecimento ao articular, de forma sistemática e fundamentada, três campos que raramente dialogam na literatura brasileira: a arteterapia de base niseana, o pensamento decolonial latino-americano e os estudos sobre parentalidade atípica. A originalidade dessa articulação reside na demonstração de que a arteterapia decolonial — herdeira do legado de Nise da Silveira, não é apenas uma técnica expressiva, mas uma prática ético-política que questiona as estruturas de poder que produzem sofrimento psíquico em populações marginalizadas. 
No plano prático, o estudo oferece subsídios para inspirações concretas para a implementação de grupos de arteterapia com familiares de crianças neurodivergentes em serviços públicos de saúde e educação. A descrição detalhada dos três encontros — com suas atividades, materiais e dinâmicas utilizados — constitui um roteiro metodológico que pode ser replicado e adaptado por outros profissionais e serviços. Ademais, a demonstração de que a arteterapia decolonial em grupo fortalece redes de apoio informais e eleva a percepção de suporte social oferece evidências que podem fundamentar políticas públicas de cuidado, alinhadas aos princípios da Política Nacional de Práticas Integrativas e Complementares (Brasil, 2017). 
Este estudo apresenta limitações que devem ser explicitadas para que seus achados sejam interpretados com a devida cautela. Quanto ao tamanho reduzido da amostra, três encontros com um grupo de familiares, limita a generalização dos resultados, embora seja compatível com a abordagem qualitativa adotada e com a profundidade analítica pretendida. A curta duração das intervenções (três meses) não permite avaliar a permanência dos efeitos ao longo do tempo. 
A arteterapia decolonial, tal como concebida e praticada neste estudo, oferece uma resposta criativa e eticamente situada ao sofrimento psíquico de familiares de crianças neurodivergentes. Ao recusar a patologização da parentalidade atípica e ao valorizar a expressão criativa como via de ressignificação, esta abordagem honra o legado de Nise da Silveira — que, décadas antes da formalização do pensamento decolonial, já praticava uma clínica que não considerava somente os diagnósticos isoladamente, mas se abria à imprevisibilidade criativa do sujeito (Mello, 2014; Silveira, 1992). 
As mães que, entre os meses de fevereiro e abril de 2026, se propuseram a participar dos encontros, pintando, desenhando e compartilhando suas experiências nos ensinaram que a parentalidade atípica não se reduz a estresse, sobrecarga e luto — é também espaço de criação, de encontro e de resistência. Se este trabalho contribuir, ainda que modestamente, para que mais espaços como esse existam — onde mães de crianças neurodivergentes possam ser cuidadas enquanto cuidam, acolhidas enquanto acolhem, e reconhecidas em sua totalidade, então terá cumprido seu propósito.Que a arteterapia decolonial siga tecendo redes de cuidado onde o colonialismo só soube impor silêncio e isolamento.

\chapter{REFERÊNCIAS}

ALVES, J. S.; GAMEIRO, A. C. P.; BIAZI, P. H. G. Estresse, depressão e ansiedade em mães de autistas: uma revisão da produção nacional. **Revista Psicopedagogia**, São Paulo, v. 39, n. 120, p. 340-352, 2022.

AMERICAN PSYCHIATRIC ASSOCIATION. **DSM-5**: manual diagnóstico e estatístico de transtornos mentais. 5. ed. Porto Alegre: Artmed, 2014. Tradução de: Diagnostic and Statistical Manual of Mental Disorders. 5. ed. Washington: APA, 2013.

AMARANTE, P. **Saúde mental e atenção psicossocial**. Rio de Janeiro: Fiocruz, 2007.

ANNOUS, N.; AL-HROUB, A.; EL ZEIN, F. A systematic review of empirical evidence on art therapy with traumatized refugee children. **Frontiers in Psychology**, v. 13, 811515, 2022. DOI: 10.3389/fpsyg.2022.811515.

ASSOCIAÇÃO BRASILEIRA DE NORMAS TÉCNICAS. **NBR 6023**: informação e documentação — referências — elaboração. Rio de Janeiro: ABNT, 2025.

ASSOCIAÇÃO BRASILEIRA DE NORMAS TÉCNICAS. **NBR 6028**: informação e documentação — resumo, resenha e recensão — apresentação. Rio de Janeiro: ABNT, 2021.

ASSOCIAÇÃO BRASILEIRA DE NORMAS TÉCNICAS. **NBR 10520**: informação e documentação — citações em documentos — apresentação. Rio de Janeiro: ABNT, 2023.

ASSIS, T. P. de; BATISTA, A. de S. M. Arteterapia como ferramenta decolonial para a saúde mental de mulheres no Quilombo de Pinhões: um olhar pela psicologia analítica. **Revista Brasileira de Promoção de Saúde**, v. 37, p. 1-15, 2024. DOI: 10.5020/18061230.2024.12345.

BOGDASHINA, O. **Sensory perceptual issues in autism and Asperger syndrome**: different sensory experiences — different perceptual worlds. 2. ed. London: Jessica Kingsley, 2016.

BRAGA DOS ANJOS, B.; ARAÚJO DE MORAIS, N. As experiências de famílias com filhos autistas: uma revisão integrativa da literatura. **Ciencias Psicológicas**, Montevideu, v. 15, n. 2, e1203, 2021. DOI: 10.22235/cp.v15i2.1203.

BRASIL. Lei nº 10.216, de 6 de abril de 2001. Dispõe sobre a proteção e os direitos das pessoas portadoras de transtornos mentais e redireciona o modelo assistencial em saúde mental. **Diário Oficial da União**, Brasília, DF, 9 abr. 2001.

BRASIL. Ministério da Educação. Secretaria de Educação Especial. **Política Nacional de Educação Especial na Perspectiva da Educação Inclusiva**. Brasília, DF: MEC/SEESP, 2008. Disponível em: [http://portal.mec.gov.br/arquivos/pdf/politicaeducespecial.pdf](http://portal.mec.gov.br/arquivos/pdf/politicaeducespecial.pdf). Acesso em: 15 abr. 2026.

BRASIL. Ministério da Saúde. Portaria nº 3.088, de 23 de dezembro de 2011. Institui a Rede de Atenção Psicossocial para pessoas com sofrimento ou transtorno mental e com necessidades decorrentes do uso de crack, álcool e outras drogas, no âmbito do SUS. **Diário Oficial da União**, Brasília, DF, 26 dez. 2011.

BRASIL. Ministério da Saúde. Conselho Nacional de Saúde. Resolução nº 466, de 12 de dezembro de 2012. Aprova as diretrizes e normas regulamentadoras de pesquisas envolvendo seres humanos. **Diário Oficial da União**, Brasília, DF, 13 jun. 2013.

BRASIL. Ministério da Saúde. Portaria nº 849, de 27 de março de 2017. Inclui a arteterapia, ayurveda, biodança, dança circular, meditação, musicoterapia, naturopatia, osteopatia, quiropraxia, reflexoterapia, reiki, shantala, terapia comunitária integrativa e yoga à Política Nacional de Práticas Integrativas e Complementares. **Diário Oficial da União**, Brasília, DF, 28 mar. 2017.

CAMPOS, R. T. O.; AMARAL, M. A. Clínica ampliada e compartilhada, a gestão democrática e redes de atenção como referenciais teórico-operacionais para a reforma psiquiátrica. **Ciência \& Saúde Coletiva**, Rio de Janeiro, v. 12, n. 4, p. 899-908, 2007. DOI: 10.1590/S1413-81232007000400014.

CARVALHAES, F. F. Clínica extramuros: decolonizando a Psicologia. **Revista Espaço Acadêmico**, Maringá, v. 19, n. 216, p. 3-13, 2019.

CATTA-PRETA, L. **Nise da Silveira**: uma psiquiatra rebelde. São Paulo: Planeta, 2021.

COMAS-D�AZ, L.; ADAMES, H. Y.; CHAVEZ-DUEÑAS, N. Y. (org.). **Decolonial psychology**: toward anticolonial theories, research, training, and practice. Washington: American Psychological Association, 2024.

CONOLLY, N. B. et al. Art group interventions for children with learning differences: a systematic review. **OTJR: Occupational Therapy Journal of Research**, v. 45, n. 2, p. 1-15, 2025. DOI: 10.1177/15394492251340378.

COSTA-ROSA, A. da. O modo psicossocial: um paradigma das práticas substitutivas ao modo asilar. **Revista Latinoamericana de Psicopatologia Fundamental**, São Paulo, v. 3, n. 1, p. 109-122, 2000. DOI: 10.1590/1415-47142000001009.

DIJKSTRA-DE NEIJS, L. et al. Parental stress and quality of life in parents of young children with autism spectrum disorder. **Child Psychiatry \& Human Development**, v. 55, n. 4, p. 1012-1025, 2024. DOI: 10.1007/s10578-024-01693-3.

DIMENSTEIN, M. et al. Estudos pós-coloniais e decoloniais na pós-graduação em psicologia no Brasil. **Psicologia Argumento**, Curitiba, v. 39, n. 105, p. 689-713, 2021. DOI: 10.7213/psicolargum.39.105.AO05.

FANON, F. **Pele negra, máscaras brancas**. Tradução Renato da Silveira. Salvador: EDUFBA, 2008.

F�VERO, M. A. B.; SANTOS, M. A. dos. Autismo infantil e estresse familiar: uma revisão sistemática da literatura. **Psicologia: Reflexão e Crítica**, Porto Alegre, v. 18, n. 3, p. 358-369, 2005. DOI: 10.1590/S0102-79722005000300010.

FOUCAULT, M. **História da loucura na idade clássica**. Tradução José Teixeira Coelho Netto. São Paulo: Perspectiva, 1978.

FRAYZE-PEREIRA, J. A. Nise da Silveira: imagens do inconsciente entre psicologia, arte e política. **Estudos Avançados**, São Paulo, v. 17, n. 49, p. 197-208, 2003. DOI: 10.1590/S0103-40142003000300012.

FREIRE, P. **Pedagogia do oprimido**. 17. ed. Rio de Janeiro: Paz e Terra, 1987.

GENTLES, S. J.; McLAUGHLIN, J.; SCHNEIDER, M. A. Stress among caregivers of autistic children: conceptual analysis and verification using two qualitative datasets. **PLOS ONE**, v. 19, n. 10, e0312391, 2024. DOI: 10.1371/journal.pone.0312391.

GOFFMAN, E. **Manicômios, prisões e conventos**. Tradução Dante Moreira Leite. São Paulo: Perspectiva, 1961.

GROSFOGUEL, R. Para uma decolonialidade dos direitos humanos: contribuições do pensamento crítico de fronteira. **Revista Direito e Práxis**, Rio de Janeiro, v. 10, n. 3, p. 1975-2012, 2019. DOI: 10.1590/2179-8966/2019/43765.

GOMBRICH, E. H. **Arte e ilusão**: um estudo da psicologia da representação pictórica. Tradução Raul de Sá Barbosa. São Paulo: Martins Fontes, 2007.

GUBA, E. G.; LINCOLN, Y. S. **Naturalistic inquiry**. Newbury Park: Sage, 1985.

GUEST, G.; BUNCE, A.; JOHNSON, L. How many interviews are enough? An experiment with data saturation and variability. **Field Methods**, v. 18, n. 1, p. 59-82, 2006. DOI: 10.1177/1525822X05279903.

HEIDEGGER, M. **A origem da obra de arte**. Tradução Maria da Conceição Costa. São Paulo: Edições 70, 2010.

JUNG, C. G. **O espírito na arte e na ciência**. 15. ed. Petrópolis: Vozes, 2000. (Obras Completas, v. 15).

KINOSHITA, R. T. Contratualidade e reabilitação psicossocial. *In*: PITTA, A. (org.). **Reabilitação psicossocial no Brasil**. São Paulo: Hucitec, 2001. p. 55-66.

MALCHIODI, C. A. (org.). **Handbook of art therapy**. 2. ed. New York: Guilford, 2012.

MAENNER, M. J. et al. Prevalence and characteristics of autism spectrum disorder among children aged 8 years — Autism and Developmental Disabilities Monitoring Network, 11 Sites, United States, 2020. **MMWR Surveillance Summaries**, Atlanta, v. 72, n. 2, p. 1-14, 2023. DOI: 10.15585/mmwr.ss7202a1.

MALDONADO-TORRES, N. On the coloniality of being: contributions to the development of a concept. **Cultural Studies**, v. 21, n. 2-3, p. 240-270, 2007. DOI: 10.1080/09502380601162548.

MANTOAN, M. T. E. **Inclusão escolar**: o que é? Por quê? Como fazer? São Paulo: Moderna, 2003.

MARETTO, L. C. Escutas de Raiz Coração: alianças afetivas e poéticas e(m) territórios e (a)travessias na vida e na arteterapia. **Revista ClimaCom - Coexistências e Cocriações**, ano 8, n. 20, 2021. Disponível em: [http://climacom.mudancasclimaticas.net.br/escutas-de-raiz/](http://climacom.mudancasclimaticas.net.br/escutas-de-raiz/). Acesso em: 15 abr. 2026.

MART�N-BARÓ, I. **Writings for a liberation psychology**. Cambridge: Harvard University Press, 1994.

McNIFF, S. **Art heals**: how creativity cures the soul. Boston: Shambhala, 2004.

MELLO, L. C. **Nise da Silveira**: caminhos de uma psiquiatra rebelde. 3. ed. Rio de Janeiro: Versal, 2014.

MERLEAU-PONTY, M. **Fenomenologia da percepção**. Tradução Carlos Alberto Ribeiro de Moura. 2. ed. São Paulo: Martins Fontes, 1999.

MIGNOLO, W. D. **The darker side of western modernity**: global futures, decolonial options. Durham: Duke University Press, 2011.

MINAYO, M. C. de S. Amostragem e saturação em pesquisa qualitativa: consensos e controvérsias. **Revista Pesquisa Qualitativa**, São Paulo, v. 5, n. 7, p. 1-12, 2017.

MOEBUS, R. L. N.; BARRETO, A. F.; MORAES, M. de M. Psicologia decolonial, contracolonial, por vir? **Estudos de Psicologia (Campinas)**, Campinas, v. 41, e230087, 2024. DOI: 10.1590/1982-0275202441e230087.

MOO, C. et al. Family-centered creative arts therapies for children with autism: a configurative systematic review. **Family Relations**, v. 73, n. 4, p. 2345-2367, 2024. DOI: 10.1111/fare.13092.

MÜLLER, A. de O. **Suturar feridas coloniais: uma leitura decolonial sobre representações de mulheres brasileiras na pintura de Marcela Cantuária**. 2022. 120 f. Dissertação (Mestrado em Educação) — Programa de Pós-Graduação em Educação, Universidade Federal do Rio Grande do Sul, Porto Alegre, 2022.

OLIVEIRA, M. de. Medicalização e decolonialidade: o dispositivo de apoio matricial em saúde mental. **Cadernos Cajuína**, v. 11, n. 3, e2215, 2026. DOI: 10.52641/cadcajv11i3.2215.

ORGANIZAÇÃO MUNDIAL DA SAÚDE. **CID-11**: Classificação Internacional de Doenças. 11. rev. Genebra: OMS, 2019. Disponível em: [https://icd.who.int/](https://icd.who.int/). Acesso em: 10 mar. 2026.

WORLD HEALTH ORGANIZATION. **World mental health report**: transforming mental health for all. Geneva: WHO, 2022. Disponível em: [https://www.who.int/publications/i/item/9789240049338](https://www.who.int/publications/i/item/9789240049338). Acesso em: 10 mar. 2026.

OLMOS-VEGA, F. et al. Decoloniality in qualitative research: a call for reflexive engagement. **Qualitative Research in Psychology**, v. 20, n. 2, p. 245-268, 2023. DOI: 10.1080/14780887.2023.2186893.

ONOCKO-CAMPOS, R. T.; FURTADO, J. P. Entre a saúde coletiva e a saúde mental: um instrumental metodológico para avaliação da rede de Centros de Atenção Psicossocial (CAPS) do Sistema Único de Saúde. **Cadernos de Saúde Pública**, Rio de Janeiro, v. 22, n. 5, p. 1053-1062, 2006. DOI: 10.1590/S0102-311X2006000500018.

PATTON, M. Q. **Qualitative research \& evaluation methods**: integrating theory and practice. 4. ed. Thousand Oaks: Sage, 2015.

PEREIRA, D. F. et al. O pensamento decolonial na psicologia brasileira: uma revisão da produção científica. **Conhecimento \& Diversidade**, Niterói, v. 14, n. 32, p. 181-193, 2022. DOI: 10.18316/rcd.v14i32.9461.

PHILIPPINI, A. **Arteterapia**: um convite ao fazer e ao sentir. Rio de Janeiro: Wak, 2020.

PINNEY, F. Neurodivergent-affirming therapeutic arts practice: becoming and companioning. **Journal of Creative Arts Therapies**, v. 17, n. 2, p. 1-17, 2022.

QUIJANO, A. Coloniality of power, eurocentrism, and Latin America. **Nepantla: Views from South**, Durham, v. 1, n. 3, p. 533-580, 2000.

RIVERA CUSICANQUI, S. **Ch'ixinakax utxiwa**: una reflexión sobre prácticas y discursos descolonizadores. Buenos Aires: Tinta Limón, 2010.

ROSAL, M. L. **Cognitive-behavioral art therapy**: from behaviorism to the third wave. New York: Routledge, 2018.

ROSE, G. **Visual methodologies**: an introduction to researching with visual materials. 4. ed. London: Sage, 2022.

RUBIN, J. A. **Introduction to art therapy**: sources and resources. 2. ed. New York: Routledge, 2010.

SANTOS, B. de S. **O fim do império cognitivo**: a afirmação das epistemologias do Sul. Belo Horizonte: Autêntica, 2019.

SCHAVERIEN, J. **The revealing image**: analytical art psychotherapy in theory and practice. London: Routledge, 1992.

SILVEIRA, N. da. **Imagens do inconsciente**. Rio de Janeiro: Alhambra, 1981.

SILVEIRA, N. da. **O mundo das imagens**. São Paulo: �tica, 1992.

SILVEIRA, N. da. **Jung: vida e obra**. 19. ed. Rio de Janeiro: Paz e Terra, 2001.

SLAYTON, S. C.; D'ARCHER, J.; KAPLAN, F. Outcome studies on the efficacy of art therapy: a review of findings. **Art Therapy: Journal of the American Art Therapy Association**, v. 27, n. 3, p. 108-118, 2010. DOI: 10.1080/07421656.2010.10129660.

SMITH, J. A.; FLOWERS, P.; LARKIN, M. **Interpretative phenomenological analysis**: theory, method and research. London: Sage, 2009.

SMITH, L. T. **Decolonizing methodologies**: research and indigenous peoples. 2. ed. London: Zed Books, 2012.

STAMPOLTZIS, A. et al. Mothers' experiences and challenges raising a child with autism spectrum disorder: a qualitative study. **Brain Sciences**, Basel, v. 11, n. 3, 309, 2021. DOI: 10.3390/brainsci11030309.

STRANG, C. E. Art therapy and neuroscience: evidence, limits, and myths. **Frontiers in Psychology**, v. 15, 1484481, 2024. DOI: 10.3389/fpsyg.2024.1484481.

URRUTIGARAY, M. H. **Arteterapia**: a transformação pessoal pelas imagens. Rio de Janeiro: Wak, 2003.

UTTLEY, L. et al. Systematic review and economic modelling of the clinical effectiveness and cost-effectiveness of art therapy among people with non-psychotic mental health disorders. **Health Technology Assessment**, Winchester, v. 19, n. 18, p. 1-120, 2015. DOI: 10.3310/hta19180.

VALLADARES-TORRES, M. M.; SOUZA, M. L. Arteterapia e sofrimento psíquico associado ao uso de drogas: revisão integrativa. **Revista Brasileira de Terapias Criativas**, Brasília, v. 1, n. 1, p. 15-28, 2024.

VAN LITH, T.; SCHOFIELD, M. J.; FENNER, P. Identifying the evidence-base for art-based practices and their potential benefit for mental health recovery: a critical review. **Disability and Rehabilitation**, v. 35, n. 16, p. 1309-1323, 2013. DOI: 10.3109/09638288.2012.732188.

VAN MANEN, M. **Researching lived experience**: human science for an action sensitive pedagogy. 2. ed. London: Routledge, 2016.

VERSITANO, S. et al. Art therapy with children and adolescents experiencing acute or severe mental health conditions: a systematic review. **Australian \& New Zealand Journal of Psychiatry**, v. 59, n. 3, p. 195-210, 2025. DOI: 10.1177/00048674251361731.

YASUI, S. **Rupturas e encontros**: desafios da reforma psiquiátrica brasileira. Rio de Janeiro: Fiocruz, 2010.

YESILKAYA, M.; MAGALLÓN-NERI, E. Parental stress related to caring for a child with autism spectrum disorder and the benefit of mindfulness-based interventions: a narrative review. **SAGE Open**, v. 14, n. 1, p. 1-15, 2024. DOI: 10.1177/21582440241235033.

YIN, R. K. **Estudo de caso**: planejamento e métodos. Tradução Daniel Grassi. 5. ed. Porto Alegre: Bookman, 2015.

ZABLOTSKY, B. et al. The economic burden of autism spectrum disorder in the United States. **Pediatrics**, v. 133, n. 5, p. e1278-e1284, 2014. DOI: 10.1542/peds.2014-0433.

ZANATTA, E. A. et al. Cotidiano de famílias que convivem com o autismo infantil. **Revista Baiana de Enfermagem**, Salvador, v. 28, n. 3, p. 271-282, 2014. DOI: 10.18471/rbe.v28i3.8177.

AITHAL, S. et al. Experiences of group dance movement psychotherapy for caregivers of children on the autism spectrum. **PLOS One**, v. 18, n. 8, e0288626, 2023. DOI: 10.1371/journal.pone.0288626.

ASSOCIAÇÃO DE ARTETERAPIA DO RIO DE JANEIRO (AARJ). **O arteterapeuta e a arteterapia decolonial**. Rio de Janeiro, 2023. Disponível em: [https://aarj.com.br/o-arteterapeuta-e-a-arteterapia-decolonial/](https://aarj.com.br/o-arteterapeuta-e-a-arteterapia-decolonial/). Acesso em: 15 abr. 2026.

AZEVEDO, E. B. et al. Arteterapia como promotora da qualidade de vida e inclusão social de profissionais e usuários. **Revista da Universidade Vale do Rio Verde**, v. 12, n. 2, p. 167-176, 2014. DOI: 10.5892/ruvrd.v12i2.1422.

BAKER, J. R. et al. Systematic review of group-based creative arts interventions in support of informal caregivers of adults: a narrative synthesis. **Ageing and Society**, v. 42, n. 6, p. 1-25, 2022. DOI: 10.1017/S0144686X2200068X.

BONAFÉ SEI, M. O grupo e as grupoterapias. In: SEI, M. B.; GONÇALVES, T. F. (Org.). **Arteterapia com grupos**: aspectos teóricos e práticos. São Paulo: Casa do Psicólogo, 2010.

CRUZ, A. C. B. O. A.; BAHIA, S. D.; SOUZA, J. C. P. A vida das mães de crianças com Transtorno de Espectro Autista: os desafios vivenciados em nossa sociedade. **Cuadernos de Educación y Desarrollo**, v. 16, n. 11, e6230, 2024. DOI: 10.55905/cuadv16n11-020.

DEBASTIANI, L. S.; ROTH, C. M. Cuidar de quem cuida: a relevância da rede de apoio na saúde mental de mães de crianças autistas. **Revista Ibero-Americana de Humanidades, Ciências e Educação**, v. 11, n. 11, p. 9708-9724, 2025. DOI: 10.51891/rease.v11i11.22744.

GRAUPEN, A.; ADRIÃO, K. G. Arteterapia promovendo respiros em tempos de incertezas. **Revista Científica de Arteterapia Cores da Vida**, v. 27, n. 16, p. 119-129, 2020.

KAÇAN, H. et al. The effect of applied group training on perceived social support, depressive symptoms, and self-compassion given to mothers with children with autism: a quasi-experimental study. **Journal of Autism and Developmental Disorders**, 2025. DOI: 10.1007/s10803-025-06963-0.

KWEK, T. Understanding community-based mental health interventions among migrant workers in Singapore. **Discover Mental Health**, v. 4, n. 41, 2024. DOI: 10.1007/s44192-024-00092-3.

LIN, B. X.; NG, V. S. B. Fostering intersectional belonging through community art therapy: a practice-based case analysis. **JoCAT**, v. 19, n. 2, 2024. Disponível em: [https://www.jocat-online.org/pp-24-lin-ng](https://www.jocat-online.org/pp-24-lin-ng). Acesso em: 15 abr. 2026.

LUNA, A. W. N. et al. Perceptions of mothers of children with autism about a support network and self-care strategies. **Revista de Enfermagem da UFPI**, v. 12, n. 1, e4284, 2023. DOI: 10.26694/reufpi.v12i1.4284.

MARINHO, E. N. M.; SILVA, A. C.; FERREIRA, M. S. Escuta psicológica sensível e acolhimento no cuidado de mães de crianças com autismo: um olhar sobre as vivências maternas. **Práxis em Saúde**, v. 3, n. 2, p. 1-8, 2025. DOI: 10.56579/prxis.v3i2.2425.

SOUSA, J. M. et al. Efetividade dos grupos terapêuticos na atenção psicossocial: análise à luz do referencial dos fatores terapêuticos. **Revista Brasileira de Enfermagem**, v. 73, e20200410, 2020. DOI: 10.1590/0034-7167-2020-0410.

SOUSA, J. M. et al. Potencialidades das intervenções grupais em Centros de Atenção Psicossocial �lcool e Drogas. **Escola Anna Nery**, v. 26, e20210294, 2021. DOI: 10.1590/2177-9465-EAN-2021-0294.

STICKLEY, T.; HUI, A. A respite thing: a qualitative study of a creative arts leisure programme for family caregivers of people with dementia. **Dementia**, v. 14, n. 3, p. 325-341, 2015. DOI: 10.1177/1471301221997290.

SUNARINGSIH, M. S.; TIATRI, S.; PATMONODEWO, S. Effectiveness of expressive art therapy group to reduce stress level in mothers of children with neurodevelopmental disorders. In: TARUMANAGARA INTERNATIONAL CONFERENCE ON THE APPLICATIONS OF SOCIAL SCIENCES AND HUMANITIES (TICASH 2019). **Proceedings** [...]. Atlantis Press, 2020. p. 606-611. DOI: 10.2991/assehr.k.200515.101.

TIRLEA, L. et al. Health, wellbeing and empowerment e-workshops for mothers of children with disabilities: a non-randomised comparison study. **Journal of Autism and Developmental Disorders**, 2024. DOI: 10.1007/s10803-024-06287-5.

WILLRICH, J. Q.; PORTELA, D. L.; CASARIN, R. Atividades de arteterapia na reabilitação de usuários da atenção psicossocial. **Revista de Enfermagem e Atenção à Saúde**, v. 7, n. 3, 2019. Disponível em: [https://seer.uftm.edu.br/revistaeletronica/index.php/enfer/article/view/3113](https://seer.uftm.edu.br/revistaeletronica/index.php/enfer/article/view/3113). Acesso em: 15 abr. 2026.

\end{document}
